\documentclass[10pt,a4paper]{amsart}
\usepackage{amsmath,amssymb,amsthm,graphicx}
\usepackage{verbatim}
\newenvironment{reference}{\comment}{\endcomment}
%\newenvironment{reference}{}{}
\newenvironment{slogan}{\comment}{\endcomment}
\newenvironment{history}{\comment}{\endcomment}
% For commutative diagrams we use Xy-pic
\usepackage[all]{xy}

% We use 2cell for 2-commutative diagrams.
\xyoption{2cell}
% Theorem environments.
%
\theoremstyle{plain}
\newtheorem{theorem}[subsection]{Theorem}
\newtheorem{proposition}[subsection]{Proposition}
\newtheorem{lemma}[subsection]{Lemma}

\theoremstyle{definition}
\newtheorem{definition}[subsection]{Definition}
\newtheorem{example}[subsection]{Example}
\newtheorem{exercise}[subsection]{Exercise}
\newtheorem{situation}[subsection]{Situation}

\theoremstyle{remark}
\newtheorem{remark}[subsection]{Remark}
\newtheorem{remarks}[subsection]{Remarks}

\numberwithin{equation}{subsection}

% Macros
%
\def\lim{\mathop{\mathrm{lim}}\nolimits}
\def\colim{\mathop{\mathrm{colim}}\nolimits}
\def\Spec{\mathop{\mathrm{Spec}}}
\def\Hom{\mathop{\mathrm{Hom}}\nolimits}
\def\Ext{\mathop{\mathrm{Ext}}\nolimits}
\def\SheafHom{\mathop{\mathcal{H}\!\mathit{om}}\nolimits}
\def\SheafExt{\mathop{\mathcal{E}\!\mathit{xt}}\nolimits}
\def\Sch{\mathit{Sch}}
\def\Mor{\mathop{\mathrm{Mor}}\nolimits}
\def\Ob{\mathop{\mathrm{Ob}}\nolimits}
\def\Sh{\mathop{\mathit{Sh}}\nolimits}
\def\NL{\mathop{N\!L}\nolimits}
\def\CH{\mathop{\mathrm{CH}}\nolimits}
\def\proetale{{pro\text{-}\acute{e}tale}}
\def\etale{{\acute{e}tale}}
\def\QCoh{\mathit{QCoh}}
\def\Ker{\mathop{\mathrm{Ker}}}
\def\Im{\mathop{\mathrm{Im}}}
\def\Coker{\mathop{\mathrm{Coker}}}
\def\Coim{\mathop{\mathrm{Coim}}}

% Boxtimes
%
\DeclareMathSymbol{\boxtimes}{\mathbin}{AMSa}{"02}

%
% Macros for moduli stacks/spaces
%
\def\QCohstack{\mathcal{QC}\!\mathit{oh}}
\def\Cohstack{\mathcal{C}\!\mathit{oh}}
\def\Spacesstack{\mathcal{S}\!\mathit{paces}}
\def\Quotfunctor{\mathrm{Quot}}
\def\Hilbfunctor{\mathrm{Hilb}}
\def\Curvesstack{\mathcal{C}\!\mathit{urves}}
\def\Polarizedstack{\mathcal{P}\!\mathit{olarized}}
\def\Complexesstack{\mathcal{C}\!\mathit{omplexes}}
% \Pic is the operator that assigns to X its picard group, usage \Pic(X)
% \Picardstack_{X/B} denotes the Picard stack of X over B
% \Picardfunctor_{X/B} denotes the Picard functor of X over B
\def\Pic{\mathop{\mathrm{Pic}}\nolimits}
\def\Picardstack{\mathcal{P}\!\mathit{ic}}
\def\Picardfunctor{\mathrm{Pic}}
\def\Deformationcategory{\mathcal{D}\!\mathit{ef}}
\usepackage{cleveref}
\crefname{originallemma}{original lemma}{original lemmas}
\Crefname{originallemma}{Original lemma}{Original lemmas}
\makeatletter
\newlabel{lemma-flat-eq}{{lemma-flat-eq}{source}{}{sourcepointer-1}{}}
\newlabel{lemma-flat-eq@cref}{{[originallemma][0][]lemma-flat-eq}{[1][0][]source}{}{}{}}
\makeatother
\title{A flat module of finite presentation is locally a direct summand of a finite free module}
\author{The Stacks Project Authors}
\date{stacks-project:modules.tex:local-snapshot-20261012; GNU FDL 1.2 or later}
\begin{document}
\maketitle
\noindent\textit{Editorial note.} All delivered mathematical text is the author's own, copied from the pinned Stacks Project chapter source sources/stacks-modules/pinned-original-source.tex; no statement or proof step was rewritten, reordered or removed.\par\medskip
\noindent\textit{Editorial note.} Preamble provenance: the Stacks Project keeps one single shared preamble for every chapter; this card takes it directly from the staged sources/stacks-modules/preamble.tex as three verbatim spans (lines 8-12, 14-18 and 154-217), so no derived file was written and nothing was added or redefined. The wrapper class and the base packages are supplied by the assembler.\par\medskip
\noindent\textit{Editorial note.} Bibliography provenance: neither the counted statement nor its proof contains a citation (this lane rescans the whole span before writing the card), so the card carries no bibliography and no entry was reconstructed.\par\medskip
\noindent\textit{Editorial note.} Selection provenance: this card was opened under the file-based control contract (CONTROL.md), with the mechanical pre-card gate run on its statement and proof interval before the card was written.\par\medskip
\noindent\textit{Allowed standard prerequisites (editorial).}\begin{itemize}
\item The local presentation of a module of finite presentation by a map between finite free modules followed by a surjection. Used at the start of the counted proof (source lines 3161-3166). Provides the two free modules whose rank controls the number of steps of the iteration.
\item The local factorisation criterion for a map followed by a surjection onto a flat module. Invoked inside the counted proof at source lines 3167-3177 and attributed to a lemma of this chapter. Allows the proof to peel off one summand of the relations at a time.
\end{itemize}\medskip
\ifdefined\hypertarget\hypertarget{sourcepointer-1}{}\fi
\noindent\textit{Source reference.} Original lemma lemma-flat-eq: sources/stacks-modules/pinned-original-source.tex, source lines 3167--3167. Navigation pointer only: the counted text refers to the Stacks item lemma-flat-eq, which belongs to another chapter of the project and is not part of this card. The pointer carries the target's own label as its number because this chapter's staged snapshot does not contain the other chapter's tag table; its content is not reproduced here.\par\medskip
\begin{lemma}
\label{lemma-flat-locally-finite-presentation}
Let $(X, \mathcal{O}_X)$ be a ringed space. Let $\mathcal{F}$ be a
flat $\mathcal{O}_X$-module of finite presentation. Then $\mathcal{F}$ is
locally a direct summand of a finite free $\mathcal{O}_X$-module.
\end{lemma}

\begin{proof}
After replacing $X$ by the members of an open covering, we may
assume there exists a presentation
$$
\mathcal{O}_X^{\oplus r} \to
\mathcal{O}_X^{\oplus n} \to \mathcal{F} \to 0
$$
Let $x \in X$. By Lemma \ref{lemma-flat-eq}
we can, after shrinking $X$ to an open
neighbourhood of $x$, assume there exists a factorization
$$
\mathcal{O}_X^{\oplus n} \to
\mathcal{O}_X^{\oplus n_1} \to \mathcal{F}
$$
such that the composition
$\mathcal{O}_X^{\oplus r} \to \mathcal{O}_X^{\oplus n} \to
\mathcal{O}_X^{\oplus n_1}$
annihilates the first summand of $\mathcal{O}_X^{\oplus r}$.
Repeating this argument $r - 1$ more times we obtain a factorization
$$
\mathcal{O}_X^{\oplus n} \to
\mathcal{O}_X^{\oplus n_r} \to \mathcal{F}
$$
such that the composition
$\mathcal{O}_X^{\oplus r} \to \mathcal{O}_X^{\oplus n}
\to \mathcal{O}_X^{\oplus n_r}$ is zero.
This means that the surjection $\mathcal{O}_X^{\oplus n_r} \to \mathcal{F}$
has a section and we win.
\end{proof}

\end{document}
